\documentclass[11pt,a4paper]{article}
\usepackage[T1]{fontenc}
\usepackage{lmodern,microtype}
\usepackage[margin=27mm]{geometry}
\usepackage{amsmath,amssymb,amsthm,booktabs,xcolor,xurl}
\usepackage[colorlinks=true,linkcolor=blue!45!black,urlcolor=blue!48!black,
  citecolor=blue!45!black,pdfauthor={Yin-Kai Yu},
  pdftitle={A Compact Lean 4 Proof of the Classical Zero-Free Region for the Riemann Zeta Function}]{hyperref}
\urlstyle{tt}
\newcommand{\lean}[1]{\nolinkurl{#1}}
\newtheorem*{maintheorem}{Theorem}
\setlength{\emergencystretch}{1.3em}
\title{A Compact Lean 4 Proof of the Classical Zero-Free Region for the Riemann Zeta Function}
\author{Yin-Kai Yu}
\date{22 November 2026}
\begin{document}
\maketitle
\begin{abstract}
We formalize an explicit de la Vall\'ee Poussin zero-free region for
the Riemann zeta function with constant \(2^{-20}\). Our local estimate
bounds the real logarithmic derivative by \(24(\log B+1)\) on a
translated disk and retains the contribution of one prescribed zero.
Factoring a zero by a differentiable divided difference halves the
growth budget. On a disk without such a zero, a normalized primitive,
Borel--Carath\'eodory and Cauchy's estimate give the same bound.
Applying it to the entire extension of \((s-1)\zeta(s)\) at three
heights makes both weighted positivity arguments available in one
calculation. Four submissions each set a new public Core~1 record.
The last uses one Lean file and
2,388 leaderboard tokens, down from 21,641 in the challenge baseline.
We supply the submitted files, replayable diffs and constants for the argument.
\end{abstract}

\section{Statement and proof architecture}
Write \(s=\sigma+it\). The formal target is the classical zero-free
region near the line \(\sigma=1\). The proof gives an explicit witness
for the existential constant in the challenge statement \cite{task}.
\begin{maintheorem}
For every \(s\in\mathbb C\) with \(s\ne 1\),
\[
   \sigma\ge 1-\frac{2^{-20}}{\log(|t|+2)}
   \quad\Longrightarrow\quad \zeta(s)\ne 0 .
\]
\end{maintheorem}
In Lean the exported declaration is \lean{AlpogeRHCore1.zeroFreeRegion}.
The witness appears in the first line of its proof. The construction
has two reusable components. First, an Abel--Mellin representation gives
both a disk growth bound and a lower bound on the height of a hypothetical
zero. Second, the local disk lemma accepts either a prescribed zero or
a zero weight. This lets all three applications share one interface. At
the middle height a zero forces a large positive reciprocal-distance term.
At the outer two heights the weight vanishes. A positive trigonometric
polynomial lets us add the three inequalities, keeping enough of the
middle term to contradict the growth bound.

\section{Abel summation and the pole at one}
\label{sec:abel-growth}

Write \(s=\sigma+it\), and let \(F\) be the entire extension of
\((s-1)\zeta(s)\), with \(F(1)=1\). The same function will be used at the
three heights \(0,t,2t\). For \(\sigma>0\), put
\[
  E(s)=\int_1^\infty \{u\}\,u^{-s-1}\,du,
  \qquad \{u\}=u-\lfloor u\rfloor .
\]
Here \(u^s=\exp(s\log u)\) for \(u>0\). We claim that
\begin{equation}
  F(s)=s\bigl(1-(s-1)E(s)\bigr)
  \qquad (\sigma>0).
  \label{eq:abel-entire}
\end{equation}

For \(\sigma>1\), the partial sums of the constant Dirichlet coefficient
\(1\) are \(\lfloor u\rfloor\). Abel summation therefore gives
\[
  \zeta(s)
  =s\int_1^\infty \lfloor u\rfloor u^{-s-1}\,du
  =s\left(\frac{1}{s-1}-E(s)\right).
\]
In the second equality we used \(\lfloor u\rfloor=u-\{u\}\) and
\(\int_1^\infty u^{-s}\,du=(s-1)^{-1}\). Both integrals in this
decomposition converge absolutely. In Lean the first equality is the
constant-coefficient specialization of \lean{LSeries_eq_mul_integral'}.

For the continuation, define \(h(u)=\mathbf{1}_{(1,\infty)}(u)\{u\}\)
on the positive real axis. Then
\(E(s)=\mathcal{M}h(-s)\), where
\(\mathcal{M}h(w)=\int_0^\infty u^{w-1}h(u)\,du\).
The measurable function \(h\) is bounded by one and vanishes in a
neighbourhood of zero. These estimates give holomorphy of the Mellin
transform for \(\Re w<0\) by
\lean{mellin_differentiableAt_of_isBigO_rpow}. Both sides of
\eqref{eq:abel-entire} are thus holomorphic on the connected half-plane
\(\sigma>0\), and the Abel identity gives equality on a neighbourhood of
\(s=2\). The identity theorem proves \eqref{eq:abel-entire} throughout the
half-plane, including at \(s=1\). These steps are formalized in
\lean{tail_smooth}, \lean{remainder_eq}, and \lean{entireFormula}.

Two consequences of the same remainder estimate will be needed. Absolute
integration gives
\begin{equation}
  |E(s)|\leq\int_1^\infty u^{-\sigma-1}\,du=\frac{1}{\sigma}.
  \label{eq:abel-remainder-bound}
\end{equation}
For \(\sigma\geq \tfrac12\), equations \eqref{eq:abel-entire}
and \eqref{eq:abel-remainder-bound} imply
\begin{equation}
  |F(s)|\leq |s|\bigl(1+2(|s|+1)\bigr)
           \leq (2|s|+2)^2 .
  \label{eq:quadratic-growth}
\end{equation}
These are \lean{mellin_bound} and \lean{growth} in the formal proof.

Suppose now that \(F(s)=0\) and \(\sigma\geq 3/4\).
Since \(F(1)=1\) and \(s\ne 0\), equation \eqref{eq:abel-entire} gives
\(E(s)=(s-1)^{-1}\). Hence \eqref{eq:abel-remainder-bound} and the triangle
inequality imply
\[
  \sigma \leq |s-1|\leq |\sigma-1|+|t|.
\]
This gives \(|t|\geq 2\sigma-1\) for \(3/4\leq\sigma\leq 1\)
and \(|t|\geq1\) for \(\sigma\geq1\). In particular,
\begin{equation}
  F(\sigma+it)=0,\quad \sigma\geq\tfrac34
  \quad\Longrightarrow\quad |t|\geq\tfrac12 .
  \label{eq:height-separation}
\end{equation}
This separation supplies the lower bound on the height used to estimate
the pole terms in the final argument.

\section{A local Landau inequality}
\label{sec:landau-disk}

Let \(D_r=\{w\in\mathbb C:|w|<r\}\). Suppose that \(f\) is entire,
\(f(0)\ne0\), every zero of \(f\) has nonpositive real part, and
\[
  |f(w)|\leq B|f(0)|\qquad (|w|\leq 1/2).
\]
In particular \(B\geq1\). Allow \(x\in\mathbb R\) either to be zero or to
satisfy \(x\leq-\Re(\rho^{-1})\) for some zero
\(\rho\in\overline D_{1/6}\). Then
\begin{equation}
  x-\Re\frac{f'(0)}{f(0)}\leq 24(\log B+1).
  \label{eq:landau-disk}
\end{equation}
The parameter \(x\) allows the same lemma to retain a specified zero at
one height and to bound the logarithmic derivative alone at the other
heights. Its two permitted cases are expressed by \lean{witness} and
the disjunction in \lean{landau}.

\subsection{Removing a zero}
\label{sec:landau-peel}

Choose a natural number \(n>2B\) and argue by induction on \(n\), allowing
\(f,B,x\) to vary. The case \(n=0\) is excluded by \(B\geq1\).
At a successor \(n=m+1\), suppose first that there is a zero \(\rho\)
as in the condition on \(x\). Since \(\rho\ne0\), the factorization
\(f(w)=(w-\rho)g(w)\) has \(g(0)\ne0\). The formal proof defines \(g\)
by the divided difference \lean{dslope}, whose value at \(\rho\) is \(f'(\rho)\).
Holomorphy of this divided difference follows from
\lean{differentiableOn_dslope}.

On \(|w|=1/2\), we have
\[
  |w-\rho|\geq \tfrac12-|\rho|\geq 2|\rho|,\qquad
  |f(0)|=|\rho|\,|g(0)|.
\]
Consequently \(|g(w)|\leq (B/2)|g(0)|\) on the circle. The maximum
modulus principle extends the estimate to its closed disk. All zeros of
\(g\) are zeros of \(f\). The induction hypothesis applies to \(g\) with
bound \(B/2\) and with \(x=0\), since \(2B<m+1\) and \(B\geq1\) give
\(2(B/2)=B<m\). Logarithmic differentiation gives
\[
 \frac{f'(0)}{f(0)}=-\rho^{-1}+\frac{g'(0)}{g(0)},\qquad
 x-\Re\frac{f'(0)}{f(0)}
   \leq -\Re\frac{g'(0)}{g(0)}
   \leq 24\bigl(\log(B/2)+1\bigr).
\]
Monotonicity of the real logarithm now gives \eqref{eq:landau-disk}.
The Lean source credits this monotonicity step after division of the
growth bound to Pi's public Core~1 submission \#96263.

\subsection{The zero-free disk}
\label{sec:landau-regular}

In the remaining branch the hypothesis on \(x\) forces \(x=0\).
Any zero \(\rho\in D_{1/6}\) would satisfy
\(-\Re(\rho^{-1})=-\Re(\rho)/|\rho|^2\geq0\), supplying a witness even
for \(x=0\). Thus \(f'/f\) is holomorphic on \(D_{1/6}\). It has a
primitive \(G\) there with \(G(0)=0\). The logarithmic derivatives of
\(f\) and \(\exp G\) agree, so their quotient is constant on the disk.
Evaluation at the origin gives \(f(w)=f(0)\exp G(w)\). The formal proof
uses \lean{isExactOn_ball} for the primitive and \lean{logDeriv_eqOn_iff}
for this identity.

Set \(A=\log B+1>0\). The original bound gives
\(\Re G(w)=\log(|f(w)|/|f(0)|)\leq\log B\) on \(D_{1/6}\).
The Borel--Carath\'eodory inequality for a function vanishing at the
centre gives
\[
  |G(w)|\leq \frac{2A|w|}{1/6-|w|}\qquad (|w|<1/6).
\]
In particular, \(|G(w)|\leq 2A\) on \(|w|=1/12\). Cauchy's derivative
estimate on that circle yields
\[
  -\Re\frac{f'(0)}{f(0)}\leq |G'(0)|
                           \leq \frac{2A}{1/12}
                           =24(\log B+1).
\]
This proves \eqref{eq:landau-disk} in the second branch. The two estimates
are implemented by \lean{borelCaratheodory_zero} and
\lean{norm_deriv_le_of_forall_mem_sphere_norm_le}.

\section{Three heights and the explicit zero-free region}
\label{sec:three-heights}

We prove that \(c=2^{-20}\) is admissible. Suppose, to the contrary, that
\(s=\beta+it\ne1\) is a zero of \(\zeta\) with
\[
  R=|t|+2,\qquad L=\log R,\qquad d=\frac{2^{-20}}{L},
  \qquad \beta\ge 1-d .
\]
All logarithms of positive real numbers are natural logarithms.
Since \(L\ge\log 2>1/2\), we have \(0<d\le 1/100\).
The nonvanishing of \(\zeta\) on \(\Re z\ge 1\), away from its pole,
gives \(\beta<1\). In particular, \(s\ne0\) and \(F(s)=0\).

Since \(\beta\ge1-d\ge 99/100\), the height separation
in \eqref{eq:height-separation} applies. More precisely, its estimate at
this zero reads
\[
    \beta\le |s-1|\le 1-\beta+|t|,
    \qquad |t|\ge 2\beta-1\ge 1-2d\ge \tfrac12 .
\]
The last, weaker lower bound suffices to control both nonreal pole terms.

\subsection{One disk estimate at three centers}

Put \(a=\beta+5d\). Then \(4d\le a-1\), so \(1<a\le 1+5d\).
For real \(v\) with \(|v|\le 2|t|\), define
\[
 z_v=a+iv,\qquad f_v(w)=F(z_v+w),\qquad
 B_v=\frac{(16R)^2}{|F(z_v)|},\qquad
 H=1+8\log 2+2L .
\]
The denominator is positive because \(a>1\).
On \(|w|\le 1/2\) we have \(\Re(z_v+w)>1/2\). Also,
\[
   |z_v+w|\le a+|v|+\tfrac12\le \tfrac{31}{20}+2|t|,
   \qquad 2|z_v+w|+2\le \tfrac{51}{10}+4|t|\le 16R .
\]
Thus \lean{ClassicalZFR.growth} gives
\(|f_v(w)|\le(16R)^2=B_v|f_v(0)|\).
If \(f_v(w)=0\), then \(\Re(z_v+w)\le1\), so \(\Re w\le1-a<0\).
These are the growth and zero-location hypotheses of \eqref{eq:landau-disk},
proved by \lean{ClassicalZFR.landau}.

Define two real quantities by
\[
   D_v=\Re\sum_{n\ge 1}\frac{\Lambda(n)}{n^{z_v}},
   \qquad P_v=\Re\frac{1}{z_v-1},
\]
where \(\Lambda\) denotes the von Mangoldt function. For \(a>1\) the
Dirichlet series converges absolutely and equals \(-\zeta'(z_v)/\zeta(z_v)\).
Logarithmic differentiation of \(F\) consequently gives
\(\Re(f_v'(0)/f_v(0))=P_v-D_v\).
This is the sign convention in \lean{ClassicalZFR.good}. With \(x=0\),
or with an admissible zero witness of weight \(x\), the disk estimate reads
\begin{equation}
   x+D_v\le 24\bigl(H-\log|F(z_v)|\bigr)+P_v .
   \label{eq:three-window}
\end{equation}
Indeed, \(1+\log B_v=H-\log|F(z_v)|\). The term \(8\log 2\) is
\(\log(16^2)\), the coefficient of \(L\) is \(2\), and \(24\) multiplies
the whole expression.

Use \eqref{eq:three-window} at \(v=0,t,2t\). The outer two applications
take \(x=0\). At \(v=t\) the translated zero is \(\rho=-5d\), since
\(z_t-5d=s\). It lies in the witness disk because \(5d\le1/20<1/6\), and
its weight is \(-\Re(\rho^{-1})=1/(5d)\). The three pole terms satisfy
\begin{align}
 P_0&=\frac{1}{a-1}\le\frac{1}{4d},&
 P_t&\le\frac{1}{|z_t-1|}\le\frac{1}{|t|},&
 P_{2t}&\le\frac{1}{|z_{2t}-1|}\le\frac{1}{2|t|}.
 \label{eq:three-poles}
\end{align}
For the last two estimates we used \(\Re u\le |u|\) and
\(|z_v-1|\ge |v|\). In particular \(4P_t+P_{2t}\le9\), since
\(|t|\ge1/2\).

\subsection{The two weighted positivity estimates}

The first estimate controls the three center values inside the logarithms
in \eqref{eq:three-window}. The Dirichlet character lemma
\lean{norm_LFunction_product_ge_one} applies to the trivial character
of modulus \(1\), real displacement \(a-1>0\), and height \(t\). It gives
\[
       1\le |\zeta(z_0)|^3|\zeta(z_t)|^4|\zeta(z_{2t})|.
\]
At each center \(|z_v-1|\ge a-1\ge4d\).
Substituting \(|\zeta(z_v)|=|F(z_v)|/|z_v-1|\) and taking real
logarithms therefore gives
\begin{equation}
 \Sigma_F:=3\log|F(z_0)|+4\log|F(z_t)|+\log|F(z_{2t})|
             \ge 8\log(4d).
 \label{eq:three-product}
\end{equation}
The exponent \(8=3+4+1\) records all the factors introduced by clearing
the pole of \(\zeta\).

The second estimate concerns \(D_v\), the Dirichlet series in the
logarithmic derivative. Absolute convergence at real part \(a>1\) allows
us to take real parts term by term and combine the three sums. For
\(\theta_n=t\log n\) the result is
\begin{align}
 3D_0+4D_t+D_{2t}
   &=\sum_{n\ge1}\Lambda(n)n^{-a}
           \bigl(3+4\cos\theta_n+\cos(2\theta_n)\bigr)\notag\\
   &=\sum_{n\ge1}2\Lambda(n)n^{-a}(1+\cos\theta_n)^2\ge0 .
 \label{eq:three-series}
\end{align}
The double-angle identity gives the square, and every coefficient
\(\Lambda(n)n^{-a}\) is nonnegative. In the formal proof
\lean{HasSum.nonneg} sums these pointwise inequalities.
Thus \eqref{eq:three-product} gives a lower bound for a sum of logarithms,
while \eqref{eq:three-series} gives the sign of a sum of derivatives.

\subsection{Closing the numerical inequality}

Multiply \eqref{eq:three-window} at \(0,t,2t\) by \(3,4,1\),
respectively, and add. By \eqref{eq:three-series} the left side is at
least \(4/(5d)\). On the right, \eqref{eq:three-poles} bounds the total
pole contribution by \(3/(4d)+9\). Hence
\[
 \frac{4}{5d}
 \le 192H-24\Sigma_F+\frac{3}{4d}+9 .
\]
Inserting \eqref{eq:three-product} and using
\(4/5-3/4=1/20\) gives
\begin{equation}
     \frac{1}{20d}\le 192\bigl(H-\log(4d)\bigr)+9 .
     \label{eq:three-final}
\end{equation}
Here \(4d=2^{-18}/L\). In particular,
\(H-\log(4d)=1+26\log 2+2L+\log L\), with \(26=8+18\).
Substituting into \eqref{eq:three-final} yields
\[
    \frac{2^{20}}{20}L
        \le 201+4992\log 2+384L+192\log L .
\]
To check the constants, use \(\log L\le L\) and \(\log 2<1\).
The last inequality would imply
\[
     (2^{20}-11520)L
       \le 4020+99840\log 2 < 103860 .
\]
But \(2^{20}-11520=1037056\) and \(L\ge\log 2>1/2\), so its left
side exceeds \(518528\). This contradiction proves
\(\zeta(s)\ne0\) throughout the claimed region for \(s\ne1\), with the
explicit constant \(2^{-20}\).

\section{What changed in the five archived proofs}
\label{sec:engineering}

The reference implementation assembles the zero-free region from named
results for zero factorization, a local logarithmic derivative estimate,
the three-height inequality, growth, a reciprocal bound, and bounded
heights. Its entry point is \lean{Formalism.ZF.zeroFreeRegion}.
The challenge assigns its 21-file dependency closure 21,641 normalized
Lean tokens. By the first snapshot studied here, the submitted proof
already used the entire function \(F(s)=(s-1)\zeta(s)\) at all three
heights. The four source comparisons below start from that implementation.

Table~\ref{tab:runs} records the public results. Comparing consecutive
archived team proofs, with v22 omitted, gives savings of 236, 168, 300
and 36 tokens. These differences measure all changes in each patch together.

\subsection{The Dirichlet series at the three centers (v21 to v23)}

In v21 the common quantity \(D(z)\) was defined as
\(-\Re(\zeta'(z)/\zeta(z))\). To prove positivity, the helper
\lean{triple} first converted three convergent Dirichlet series into this
form. In v23, \(D(z)\) is defined directly as the real part of the
von Mangoldt Dirichlet series. The library result
\lean{LSeries_vonMangoldt_eq_deriv_riemannZeta_div}
identifies that series with \(-\zeta'/\zeta\) for \(\Re z>1\).
The conversion now occurs once, in the identity relating \(D\) to
\(F'/F\). At the point of summation,
\lean{LSeriesSummable_vonMangoldt} gives convergence at every required
height. With \(S(v)\) the real \lean{HasSum} statement at \(a+iv\),
the submitted combination reads
\begin{verbatim}
(((S 0).mul_left 3).add ((S s.im).mul_left 4) |>.add (S (2*s.im)))
\end{verbatim}
\lean{HasSum.nonneg} transfers termwise nonnegativity to this sum.
Expanding the cosine of twice the angle turns the trigonometric factor
into \(2(1+\cos\theta)^2\).

The same edit moved the uses of \lean{zeroHeight} and \lean{poleTerm_le}
into the contradiction argument. It also simplified the passage from
the modulus-one character product
\lean{DirichletCharacter.norm_LFunction_product_ge_one} to a product of
values of \(F\). Given the horizontal gap \(d\leq a-1\), the single
inequality \(d\leq |a+iv-1|\) holds for every \(v\). Rewriting
\(\zeta(z)=(z-1)^{-1}F(z)\) in the product exposes all three
denominators to the same order-preserving rewrite. The earlier script
constructed and applied a separate bound on \(|\zeta(a+iv)|\) at each
height. See \texttt{diffs/96267-to-96281.diff}.

\subsection{A derivative bound from one circle (v23 to v24)}

In the zero-free branch of the disk argument, let \(g\) denote the
primitive of \(f'/f\), normalized by \(g(0)=0\), and put
\(A=\log B+1\). Both proofs establish \(\Re g\leq A\) on the disk of
radius \(1/6\). In v23, \lean{HasDerivAt.le_of_lip'} required an
eventual estimate \(|g(w)-g(0)|\leq C|w|\) at zero. The proof therefore
introduced a neighborhood filter and kept the variable \(|w|\) throughout
the estimate. In v24, \lean{Complex.borelCaratheodory_zero} supplies
\[
 |g(w)|\leq \frac{2A|w|}{1/6-|w|}.
\]
Evaluating on \(|w|=1/12\) gives \(|g(w)|\leq 2A\).
The Cauchy estimate
\lean{Complex.norm_deriv_le_of_forall_mem_sphere_norm_le} takes a bound
\(C\) on a sphere of radius \(r\), together with differentiability inside
and continuity on the closure, and returns \(|g'(0)|\leq C/r\).
Here \(C/r=24A\). Once the boundary estimate is in the hypothesis
\lean{H}, the archived script closes its numerical step with
\begin{verbatim}
        simp at hz
        norm_num [hz] at H;linarith
\end{verbatim}

The interface between the disk estimate and the main theorem changed
as well. The helper \lean{good} now returns a complex equality for \(F'/F\)
alongside nonvanishing. The main proof instantiates a single local
function at \(0,t,2t\) and rewrites all three resulting inequalities
together. This removes the separately stated \lean{scaledWindow} and the
real abbreviation \(D\). In the zero-removal branch, monotonicity of
\(\log\) gives \(\log(B/2)\leq\log B\) directly. The source credits
this last observation to Pi's public submission \#96263.
See \texttt{diffs/96281-to-96285.diff}.

\subsection{The induction parameter and the growth bound (v24 to v25)}

Removing a zero \(\rho\) within radius \(1/6\) gives a factorization
\(f(w)=(w-\rho)g(w)\), where the relative growth bound for \(g\) on
the disk of radius \(1/2\) is \(B/2\). In v24 the induction starts with
\(B<2^n\). In v25, \lean{exists_nat_gt} supplies \(n>2B\).
At a successor step \(2B<n+1\) and \(B\geq 1\) imply \(B<n\).
This is precisely the induction premise \(2(B/2)<n\) for the new
function. Linear arithmetic now suffices for the decrease.

The two versions use the same divided difference \lean{dslope}.
The library lemma \lean{differentiableOn_dslope} transfers
differentiability across the removed zero, and
\lean{sub_smul_dslope_of_zero} identifies the product with \(f\)
pointwise. In v25, function extensionality turns the latter into an
equality of functions. Destructuring the existence statement with
\lean{rfl} replaces \(f\) by its product in every hypothesis and in the
goal. In particular, \lean{logDeriv_mul} now applies directly to the
goal. Its hypotheses are nonvanishing and differentiability of both
factors at zero. The maximum-modulus argument still uses
\lean{Complex.norm_le_of_forall_mem_frontier_norm_le} to pass from the
boundary estimate to the closed disk.

Independently, the growth theorem changed its mathematical statement.
The v24 bound was \(|F(z)|\leq R^8\) under \(R\geq2\),
\(\Re z\geq1/2\) and \(|z|\leq4R\). The v25 lemma gives
\(|F(z)|\leq(2|z|+2)^2\) under the sole hypothesis \(\Re z\geq1/2\).
The disk proof specializes it to \((16R)^2\) at all three centers.
The Mellin remainder is now defined by its integral, the form needed
for the norm bound. A function equality identifies it with the Mellin
transform at \(-z\) when differentiability is needed.
\lean{Complex.isOpen_re_gt} supplies the open half-planes used in
analytic continuation. See \texttt{diffs/96285-to-96293.diff}.

\subsection{Two algebraic end steps (v25 to v26)}

In the Abel identity, v26 asks directly for the factored expression
\(\zeta(z)=z((z-1)^{-1}-E(z))\), with \(E(z)\) the remainder integral.
The partial summation formula \lean{LSeries_eq_mul_integral'}
already has \(z\) outside its integral. The new expression therefore
matches the conclusion of integral subtraction without redistributing
this factor.

In the branch of the induction that removes a zero, the recursive inequality
\lean{ih0} contains \(\log(B/2)\) on its right-hand side. From \(1\leq B\),
the proof uses \(B/2\leq B\) inside that inequality, then closes the
linear goal. The two lines in v26 are
\begin{verbatim}
      grw [half_le_self (zero_le_one.trans hB)] at ih0
      linarith
\end{verbatim}
Finally, the summand with \(\theta=t\log n\) and
\(w=\Lambda(n)\exp(-a\log n)\) reduces to
\(w(3+4\cos\theta+\cos(2\theta))\).
Previously a \lean{calc} block proved \(2w(1+\cos\theta)^2\geq0\) and
then proved the polynomial identity by \lean{ring}. In v26,
\lean{linear_combination} multiplies the inequality for the square by
\(2w\) and checks the resulting identity in one step. These changes,
together with the smaller elaboration edits in
\texttt{diffs/96293-to-96295.diff}, account for the final 36-token drop.


\section{What the competition records measure}
\label{sec:records}
The challenge measures the source needed to prove a fixed Lean type.
Levent Alp\"oge's reference closure for Core~1 contains 21 Lean files,
74,010 characters and 21,641 normalized tokens \cite{task}. Table~\ref{tab:runs}
separates the public source metric from the final score returned by the
Playground evaluator. The previous best in each row belongs to the same
pillar immediately before that attempt. Competing entries between our
submissions account for the difference between the two columns of savings.

\begin{table}[htbp]
\centering\small
\begin{tabular}{@{}rrrrrr@{}}
\toprule
Attempt & Board tokens & Previous best & Record gain & Own reduction & Final score\\
\midrule
96267 & 3,128 & 3,099 & --- & --- & 64.468460\\
96281 & 2,892 & 3,024 & 132 & 236 & 64.892842\\
96285 & 2,724 & 2,762 & 38  & 168 & 65.226270\\
96293 & 2,424 & 2,686 & 262 & 300 & 65.798238\\
96295 & 2,388 & 2,424 & 36  & 36  & 65.873483\\
\bottomrule
\end{tabular}
\caption{Selected scored submissions, in chronological order. Own reduction
compares adjacent rows. A dash denotes the comparison entry or an inapplicable gain.}
\label{tab:runs}
\end{table}

At the archived 22 November 2026 snapshot (06:32:56, UTC+8), attempt
\#96295 held the Core~1 best with an \(88.97\%\) reduction from the
reference closure. Across all five pillars, the public curve advanced at
\#96281, \#96285 and \#96293. The final 36-token improvement was also marked
as a retained breakthrough. The five-pillar percentage rounded to \(41.61\%\)
on both sides of that last point \cite{records}. The public archive preserves
the unrounded final scores as reported, to their available precision.

\section{Artifacts and reproduction}
The research archive \cite{archive} includes the five submitted
\lean{Solution.lean} files in \lean{snapshots/}, four full unified patches
in \lean{diffs/}, and the record data in \lean{records/results.json}.
The final source is at the project root. The archive verifier applies
each patch in turn and checks every reconstructed source against its
SHA-256 digest and leaderboard token count. The sources were extracted
from the submission bundles identified by the CLI receipts.

For mathematical replay, the package uses Lean 4.34.1 with Mathlib
revision \lean{d13f23b723b8a846827a245b89c10fc7d3f11612}.
The standalone audit imports the compiled solution, checks the full
challenge type, and prints the axioms of the exported theorem. Those
axioms are \lean{propext}, \lean{Classical.choice} and \lean{Quot.sound}.
The repository includes commands for both the Lean build and the
separate \lean{leanchecker} replay.

\paragraph{Credits.}
The zero-free region and the three-height positivity argument are classical.
The competition supplies the formal target and uses Alp\"oge's source
release for its baseline \cite{alpoge}.
Mathlib supplies the analytic and number-theoretic lemmas used throughout
\cite{mathlib,bc}. Pi's public submission \#96263 documents the monotone
bound after halving the growth budget in the factor-removal branch
\cite{pi}. Yin-Kai Yu's entries were produced with GPT-6.2 Astra using
ChatGPT Desktop. The numbered diffs identify the changes in this series.

\begin{thebibliography}{9}
\bibitem{task} Agentic Science Challenge.
\href{https://play.bohrium.com/challenges/lean-4-verification-rh-core-pillar-1-5-classical-z-c7f72729}{Core Pillar 1 of 5, classical zero-free region for the Riemann zeta function}.
Task statement, reference closure and evaluation protocol.
\bibitem{alpoge} L. Alp\"oge.
\href{https://alpo.ge/remark.pdf}{Author's note accompanying the Lean source release},
footnote~1. The reference declaration in the
\href{https://alpo.ge/lean.tar.gz.part01}{multivolume source archive}
is \lean{Formalism.ZF.zeroFreeRegion} in \lean{src/Formalism/ZeroFreeRegion/FinalZFR.lean}.
\bibitem{mathlib} The mathlib community.
\href{https://github.com/leanprover-community/mathlib4/tree/d13f23b723b8a846827a245b89c10fc7d3f11612}{Mathlib for Lean 4, revision d13f23b7}.
\bibitem{bc} M. Radziwi\l{}\l{}.
\href{https://github.com/leanprover-community/mathlib4/blob/d13f23b723b8a846827a245b89c10fc7d3f11612/Mathlib/Analysis/Complex/BorelCaratheodory.lean}{Borel--Carath\'eodory theorem}. In Mathlib, pinned revision.
\bibitem{pi} Pi.
\href{https://riemann.verifytask.cn/api/attempt/96263}{Public Core~1 submission \#96263}.
\bibitem{records} Agentic Science Challenge.
\href{https://riemann.verifytask.cn/}{Proof Compression Progress Curve} and
\href{https://riemann.verifytask.cn/api/attempt/96295}{public record for \#96295}.
\bibitem{archive} Y.-K. Yu.
\href{https://github.com/YinkaiYu/zeta-zero-free-lean}{A compact Lean 4 proof of the classical zero-free region}. Source, diffs, data and paper, 2026.
\end{thebibliography}
\end{document}
